% =============================================================
% CHAPTER 5 — EXPERIMENTAL SETUP
% Purpose: every concrete detail needed to REPRODUCE your experiments. If
% Chapter 4 is "what and why", this is "exactly how, with which numbers". A
% careful reader should be able to rerun your study from this chapter.
% =============================================================
\chapter{Experimental Setup}
\label{chap:setup}

This chapter records the concrete configuration of the experiments:
the datasets and how they were sampled, the multimodal large language models
(MLLMs) evaluated, the prompts used to elicit a verdict, the decoding settings,
and the hardware and software environment. Where \Cref{chap:methodology}
explains \emph{what} was done and \emph{why}, this chapter states \emph{exactly
how}, so that the study can be reproduced from the open configuration files.

The experiments are organised as a funnel of zero-shot stages followed by a
separate group of few-shot in-context-learning (ICL) studies. \emph{Stage~1} screens
all \num{9} models against \num{28} label-wording variants; \emph{Stage~2}
sweeps prompt structure and persona for the models carried forward; \emph{Stage~3}
adds explicit chain-of-thought (CoT) reasoning; and the \emph{few-shot} studies
prepends labelled demonstration images. The two screening stages run on a small,
class-balanced subset of each dataset to keep the large prompt grid affordable,
and the winning configurations are then re-evaluated on the \emph{full}
datasets. Every run is one cell of a Task~$\times$~Models~$\times$~Datasets~$\times$~Prompts
matrix (\Cref{chap:methodology}); the concrete values of each axis are given
below.

\section{Datasets}
\label{sec:datasets}

Two face-swap datasets are used, \textbf{FOWS} (\emph{Face Occlusion With
Swapping})~\cite{fows} and \textbf{GOTCHA}~\cite{gotcha}. Both are collections of
video frames of a single subject, organised into an \emph{occluded} and a
\emph{non-occluded} condition, and both pair, for every subject and challenge, an
original (\emph{real}) recording with one or more face-swapped (\emph{fake})
versions produced by different generators. The occlusions are produced by
\emph{challenges}: the subject is asked to move a hand or a supplied object in
front of their face, which is known to induce visible artefacts in face-swapped
frames. Because no model is trained or
fine-tuned, the datasets' own train/test partition is irrelevant to the zero-shot
evaluation; it matters only for the few-shot studies, whose demonstration images
are drawn exclusively from the training split (\Cref{sec:prompts}).

FOWS was collected by Ziglio et al.~\cite{fows} from seven volunteers performing
hand- and object-occlusion challenges; the version used here provides one real
source and four face-swap generators (SimSwap, GHOST, FaceDancer, and
DeepFaceLab, \texttt{dlc}), yielding a strongly fake-skewed class balance (four
fake frames per real frame). GOTCHA~\cite{gotcha} is used through the subset and
protocol defined by Ziglio et al.~\cite{fows}: only the hand- and object-occlusion
challenges are retained, with two swap generators, DeepFaceLab (\texttt{DFL}) and
FSGAN, giving a mildly real-skewed balance. \Cref{tab:datasets} reports the exact
composition of the full sets as evaluated; the class imbalance in both (in
opposite directions) is the reason the evaluation in \Cref{chap:results} reports
balanced metrics rather than plain accuracy.

\begin{table}[t]
	\centering
	\caption{Full datasets used in the evaluation, by class, occlusion
	condition, and generator. Counts are frames.}
	\label{tab:datasets}
	\begin{tabular}{llrr}
		\toprule
		 & & \textbf{FOWS} & \textbf{GOTCHA} \\
		\midrule
		\multirow{2}{*}{Class}
		 & Real (original)     & \num{8400}  & \num{83771} \\
		 & Fake (swapped)      & \num{33600} & \num{76264} \\
		\midrule
		\multirow{2}{*}{Occlusion}
		 & Occluded            & \num{21000} & \num{101315} \\
		 & Non-occluded        & \num{21000} & \num{58720} \\
		\midrule
		\multirow{6}{*}{Generator}
		 & Original (real)     & \num{8400}  & \num{83771} \\
		 & SimSwap             & \num{8400}  & --          \\
		 & GHOST               & \num{8400}  & --          \\
		 & FaceDancer          & \num{8400}  & --          \\
		 & DeepFaceLab         & \num{8400}  & \num{37875} \\
		 & FSGAN               & --          & \num{38389} \\
		\midrule
		\textbf{Total} & & \num{42000} & \num{160035} \\
		\bottomrule
	\end{tabular}
\end{table}

\paragraph{Balanced screening subset.}
Running the large Stage-1 and Stage-2 prompt grids over the full datasets would
be prohibitively expensive, so the two screening stages use a fixed, class- and
condition-balanced subset of \num{1344} frames per dataset. The subset is
balanced across the label (\num{672} real / \num{672} fake) and the occlusion
condition (\num{672} occluded / \num{672} non-occluded), and is stratified by
generator: for FOWS, \num{672} real frames and \num{168} frames from each of the
four fake generators; for GOTCHA, \num{672} real frames and \num{336} frames
from each of the two fake generators. The same frozen subset is reused across
all screening runs so that per-model prompt comparisons are made on identical
images. The winning configuration of each stage is then re-evaluated on the full
datasets of \Cref{tab:datasets}.

\section{Models}
\label{sec:models}

Nine open-weight MLLMs were evaluated, spanning dense and mixture-of-experts
(MoE) architectures and a range of parameter counts. They are listed in
\Cref{tab:models} together with the exact Hugging Face checkpoint used to
identify them unambiguously. All models are run through the same vLLM backend
(\Cref{sec:hwenv}) in their instruction-tuned/chat form and are used strictly
zero-shot: no fine-tuning, and no LoRA training on the task.

For models that expose an optional ``thinking'' mode, that mode is disabled by
default (\texttt{enable\_thinking = false}, via a custom chat template where
required) so that the zero-shot verdict is a direct answer; native thinking is
studied separately and deliberately in Stage~3. Two models need model-specific
handling: \texttt{Phi-4-multimodal-instruct} is served with its vision LoRA
adapter and a batch size of~1, and \texttt{Llama-4-Scout} is the
\texttt{w4a16}-quantised release. Context lengths are left at the model default
except where a checkpoint required an explicit cap to fit memory (recorded in the
configuration files).

\begin{table}[t]
	\centering
	\caption{Multimodal LLMs evaluated. All are run zero-shot through vLLM with
	greedy decoding; ``A3B'' denotes a mixture-of-experts model with $\approx$3B
	active parameters.}
	\label{tab:models}
	\begin{tabular}{lll}
		\toprule
		Model (alias) & Checkpoint & Type \\
		\midrule
		Qwen3.6-35B-A3B      & \texttt{Qwen/Qwen3.6-35B-A3B}                     & MoE (A3B) \\
		Qwen3.6-27B          & \texttt{Qwen/Qwen3.6-27B}                         & Dense \\
		gemma-4-31B-it       & \texttt{google/gemma-4-31B-it}                    & Dense \\
		Kimi-VL-A3B-Instruct & \texttt{moonshotai/Kimi-VL-A3B-Instruct}          & MoE (A3B) \\
		InternVL3.5-30B-A3B  & \texttt{OpenGVLab/InternVL3\_5-30B-A3B-HF}         & MoE (A3B) \\
		Nemotron-Nano-12B-VL & \texttt{nvidia/NVIDIA-Nemotron-Nano-12B-v2-VL}     & Dense \\
		Llama-4-Scout-17B    & \texttt{RedHatAI/Llama-4-Scout-17B-16E \dots w4a16} & MoE (16E), w4a16 \\
		Phi-4-multimodal     & \texttt{microsoft/Phi-4-multimodal-instruct}      & Dense + vision LoRA \\
		GLM-4.6V-Flash       & \texttt{zai-org/GLM-4.6V-Flash}                    & Dense \\
		\bottomrule
	\end{tabular}
\end{table}

\section{Prompts}
\label{sec:prompts}

The staged prompt design and the reasoning behind it are described in
\Cref{sec:prompting}; this section records the concrete variants used at each
stage. Every prompt asks for a single-image verdict and constrains the answer to
two words (one standing for \emph{real} and one for \emph{fake}), which the
response parser maps back to a label (\Cref{sec:metrics}). Full prompt texts are
given in \Cref{app:prompts}; the naming introduced here is reused in
\Cref{chap:results}.

\paragraph{Stage 1: label wording.}
Stage~1 fixes a single-turn, direct-answer template and varies only the two
words that denote the classes. Four ``real'' words \{\emph{authentic}, \emph{bona
fide}, \emph{original}, \emph{real}\} are crossed with seven ``fake'' words
\{\emph{fake}, \emph{deepfake}, \emph{generated}, \emph{faceswapped},
\emph{manipulated}, \emph{synthetic}, \emph{modified}\}, giving \num{28} wording
variants. All \num{9} models are run over all \num{28} variants on the balanced
subset of both datasets. The (real, fake) word pair selected for each
carried-forward model, by the criterion of \Cref{sec:protocol}, is fixed
for the remaining stages (\Cref{tab:stage2words}).

\begin{table}[t]
	\centering
	\caption{Stage-1 label wording carried forward per model into Stage~2. Each
	pair is selected on the balanced subset; selection results are reported in
	\Cref{chap:results}.}
	\label{tab:stage2words}
	\begin{tabular}{lll}
		\toprule
		Model & ``real'' word & ``fake'' word \\
		\midrule
		Qwen3.6-35B-A3B & authentic & faceswapped \\
		Qwen3.6-27B     & bona fide & faceswapped \\
		gemma-4-31B-it  & bona fide & generated \\
		Kimi-VL-A3B     & original  & synthetic \\
		Llama-4-Scout   & authentic & fake \\
		\bottomrule
	\end{tabular}
\end{table}

\paragraph{Stage 2: prompt structure and persona.}
With the wording fixed, Stage~2 varies the \emph{structure} of the prompt. Each
carried-forward model is swept over \num{34} structures: four base templates
that place either a short task description or a definition of ``face swap'' in
either the system or the user message, plus five persona templates each
instantiated with six analyst personas (deepfake specialist, forensic examiner,
VFX artist, biometric analyst, computer-vision researcher, and content-integrity
moderator). This sweep runs on the balanced subset; the single best structure
per model is then re-evaluated on the full datasets.

\paragraph{Stage 3: chain-of-thought.}
Stage~3 takes each model's winning Stage-2 prompt and adds an explicit
instruction to reason before answering (``first briefly explain your reasoning,
then answer with exactly one phrase''). A complementary variant instead enables
the model's \emph{native} thinking mode on the same prompt, over a small fixed
set of frames, so that the raw reasoning trace can be inspected qualitatively.
Stage~3 is run on both the balanced subset and the full datasets.

\paragraph{Few-shot in-context learning.}
The few-shot studies prepend $n=10$ labelled demonstration images to the test
image. Demonstrations are drawn only from the training split and organised into
four frozen \emph{pools} (one per dataset $\times$ occlusion condition), each
holding ten real and ten fake exemplars. Demonstrations are varied by
\emph{composition} (e.g.\ all-fake, all-real, interleaved), by presentation
\emph{method} (a single user message vs.\ one chat turn per demonstration), and
by \emph{order}; cross-dataset pairings (e.g.\ FOWS demonstrations with a GOTCHA
test image) are included (see \Cref{sec:prompting} for the rationale). Few-shot
results are compared to the zero-shot Stage-2 baseline \emph{per image} (a
row-level join on image identity) so that the measured difference reflects the
demonstrations alone.

\section{Decoding and Hyperparameters}
\label{sec:decoding}

All experiments use \textbf{greedy, deterministic decoding}: the sampling
temperature is fixed at \num{0.0} and the vLLM engine seed at~\num{0}. Combined
with the fact that the models are frozen, this makes each (model, prompt, image)
verdict reproducible run-to-run. The generation budget is \num{512} tokens for
the direct-answer stages (Stages~1--2), raised to \num{1024} for the
chain-of-thought stage (Stage~3), and higher still for the native-thinking
variant, whose traces run long and would otherwise be truncated before the final
verdict.

Batch size affects only throughput, not the outputs. Zero-shot stages use a
batch size of \num{150} (\num{1} for \texttt{Phi-4-multimodal}, which is served
with its vision LoRA), and the few-shot studies use \num{32}, reflecting the
$\sim\!11\times$ larger prompt of carrying eleven images per request. Exactly one
model is loaded per operating-system subprocess, following the memory-management
strategy of \Cref{sec:pipeline}.

\section{Hardware and Software Environment}
\label{sec:hwenv}

All inference was run on a single \textbf{NVIDIA A100 80\,GB} GPU
(\texttt{tensor\_parallel\_size = 1}) in \texttt{bfloat16}. Serving used
\textbf{vLLM~0.20.2} (offline batched inference, prefix caching enabled) on a
\textbf{CUDA~13.0} PyTorch build, under \textbf{Python~3.12}; dependencies are
pinned and reproduced with \texttt{uv} from the project lockfile. vLLM's offline
batched serving~\cite{vllm} is what makes evaluating tens of thousands of frames
across nine models tractable on a single GPU. All checkpoints and their per-model
settings (chat templates, context caps, quantisation, LoRA) are stored as
version-controlled TOML configuration files alongside the study definitions.

\section{Evaluation Metrics and Reproducibility}
\label{sec:metrics}

\paragraph{From free text to a label.}
Each model returns free text, which is reduced to a verdict by a
\emph{single-phrase} parser: it scans the output for the prompt's declared real
and fake words and takes the first (or last) match as the predicted label. An
output containing neither word is recorded as a \emph{parse failure}, kept
distinct from a wrong answer (the rationale is given in \Cref{sec:parsing}). The
ground-truth label of each frame is determined by its generator folder
(original~$=$~real, otherwise~$=$~fake).

\paragraph{Metrics.}
Two metric roles are distinguished (see \Cref{sec:protocol}). Winner
\emph{selection} inside the funnel uses the Matthews Correlation Coefficient
(MCC). Final \emph{reporting} in \Cref{chap:results} uses accuracy, precision,
recall, and F1, plus \textbf{balanced accuracy} and macro-averaged F1 (all
defined in \Cref{chap:background}), which are the primary figures of merit given
the class imbalance of both datasets (\Cref{tab:datasets}). Parse-failure rate is
reported alongside. Metrics are computed both in aggregate and broken down by
dataset, occlusion condition, and generator, so that generalisation behaviour
across manipulation types is visible.

\paragraph{Reproducibility.}
Determinism is controlled by greedy decoding (temperature~\num{0}) and a fixed
engine seed (\Cref{sec:decoding}). Raw outputs are written as append-only JSONL,
keyed by image path, which both makes runs resumable and preserves every raw
generation for re-parsing or error analysis. The pipeline, study definitions,
prompts, and model configurations are all version-controlled, so any reported
number can be traced back to the exact configuration that produced it.
